\chapter{Conclusion}

\section{Achieved Results}
The integration of the SmartEdu system has successfully demonstrated a working, agent-driven educational platform. By focusing on practical functionality and stable outputs, the system has achieved the following results:

\begin{itemize}
    \item {Unified Knowledge Graph:} The system automatically builds a connected and non-fragmented Knowledge Graph. Concepts remain consistent throughout the course, providing a reliable database for all agent tasks.
    \item {Consistent AI Operations:} Guided by state-based harnessing, the multi-agent system runs stably. It reliably retrieves data, maps out learning roadmaps, and identifies backbone concepts without hallucinating or drifting off-topic.
    \item {Scalable System Design:} The backend architecture is modular and separates state management from AI processing, making it ready to scale into a microservices environment.
    \item {Multimodal Input Processing:} The platform successfully preserves the structural layouts and visual references of the original materials, allowing agents to "read" and utilize this rich context for better answers.
    \item {Reliable Error Handling:} The system includes a robust fallback mechanism. When specific context is missing, the agents safely switch to exploring the broader graph instead of generating fake information.
\end{itemize}

\section{System Limitations}
Despite its functional success, the current version of the system has several practical limitations that affect its speed and long-term usability:

\begin{itemize}
    \item {Unclear Course-Level Paths:} Defining prerequisites for an entire course is still difficult. The exact threshold for deciding when to merge overlapping topics remains undecided, which can cause confusion in the macro-roadmap.
    \item {Lack of Self-Improvement:} The system's intelligence comes from static, specialized prompts (the harness), not the models themselves. The agents cannot learn from their own mistakes, meaning the system is not yet truly adaptive.
    \item {High Latency:} Processing time is a major bottleneck. A simple data retrieval takes up to 30 seconds, while complex tasks like generating a roadmap or a teaching response can take up to 60 seconds.
    \item {Limited Long-Term Tracking:} The context from which Teaching Assistant assesses situation are still tied to the present chat session. 
    It currently lacks a reliable way to track and evaluate a student's progress over a long period.
    \item {Incomplete Application:} The frontend interface is not developed, and the application is still not hosted, limiting how users can visually interact 
    with the system and their roadmaps.
\end{itemize}

\section{Future Plan}
To prepare the SmartEdu system for final deployment and defense, the remaining work is scheduled into an 8-week execution plan. 
This focuses on visual interfaces, speed optimization, and final testing:

\begin{enumerate}
    \item \textbf{Weeks 1: Completion of Agentic Design on a Macro Scale} Improve agentic module performance for more challenging cases. 
    Propose reliable metrics.
    \item \textbf{Weeks 2: Interactive UI/UX Design.} Build the frontend interface and and test system deployments. Provide addtional features:
    \begin{itemize}
        \item Let students visually navigate the Knowledge Graph, view their roadmaps dynamically, and chat easily with the system.
        \item Streamline output to reduce response latency, display original PDF file with highlighter for deeper knowledge comprehension
    \end{itemize}
    
    \item \textbf{Weeks 3--4: Student Tracking and Exercises.} Add interactive exercises and a reliable scoring and evaluation system for student. Using these academic records, the system will apply clustering algorithms to group students with similar learning patterns for better long-term tracking.

    \item \textbf{Weeks 5--6: Optimization of Agent components on a Micro Scale } Improve Agentic Memory, Algorithm construction and speed up inference. 
    This step aims to cut down the current 30--60, while still improving quality.
    
    
    \item \textbf{Weeks 7--8: System Benchmarking and Final Report.} Run practical tests to evaluate the TA's accuracy and its impact on students. 
    Finally, summarize all results into the graduation thesis and prepare the visual presentation for the defense.
\end{enumerate}